%% Согласно ГОСТ Р 7.0.11-2011:
%% 5.3.3 В заключении диссертации излагают итоги выполненного исследования, рекомендации, перспективы дальнейшей разработки темы.
%% 9.2.3 В заключении автореферата диссертации излагают итоги данного исследования, рекомендации и перспективы дальнейшей разработки темы.
%% Поэтому эта часть общая и загружается из одного файла в автореферат и в диссертацию.
%%
%% ВНИМАНИЕ: файл подключается и диссертацией (Dissertation/conclusion.tex), и авторефератом
%% (Synopsis/content.tex). Ссылки \cref/\ref здесь использовать нельзя: метки глав и разделов
%% существуют только в диссертации, в автореферате они дадут неразрешённые ссылки.
%% Указания на главы и разделы приводятся текстом.
\begin{enumerate}
	\item Выполнен анализ моделей окружающего пространства БТС, алгоритмов их построения и способов комплексирования разнородных данных. На его основе обоснован выбор карты проходимости пространства в качестве модели, описывающей проходимость окружающего БТС пространства, байесовского вывода с использованием обратной модели наблюдения в качестве алгоритма её построения, а также коэффициента Жаккара и метрики PFC-MSE в качестве метрик оценки её качества. Сформулированы требования к разрабатываемым модели и алгоритмам её построения.
	\item Предложен алгоритм оценки положения окружающих ТС на основе комплексирования визуальных и лидарных данных, выделяющий точки поверхности движения оценкой её формы регрессией на гауссовских процессах. На реальных данных алгоритм показал средний коэффициент Жаккара $0,439$ против $0,399$ у лучшего лидарного и $0,337$ у визуального алгоритмов, не использующих обучаемых преобразований в вид сверху. Алгоритм сохраняет превосходство во всём диапазоне рассмотренных дальностей.
	\item Предложен алгоритм проецирования семантической сегментации проезжей части с изображения на оценённую по лидарным данным поверхность движения, а также способ его объединения с базовым алгоритмом построения карты проходимости пространства, использующий конкурирующее комплексирование. Результаты экспериментов показали, что полученное объединение алгоритмов снижает среднеквадратичную ошибку стоимости поиска пути (PathFinding Cost Mean Squared Error, PFC-MSE) с $1019,3$ до $924,4$ в пяти рассмотренных сценариях, а среднеквадратичную ошибку относительно эталонной карты -- с $0,231$ до $0,218$. Проецирование той же сегментации на фиксированную плоскость увеличивает PFC-MSE до $1059,4$. 
	\item Разработан программный комплекс построения карты проходимости пространства на основе комплексирования разнородных данных от датчиков, установленных на БТС, и их обработчиков. Архитектура программного комплекса позволяет осуществлять параллельную подготовку фрагментов карты по показаниям независимых источников данных, оставляя последовательным только этап добавления фрагментов в общую карту.
	\item Реализованы три последовательные модификации структуры данных для представления карты проходимости пространства, направленные на снижение вычислительной сложности производимых операций. Тестирование на синтетических данных, моделирующих одномерную карту проходимости пространства, показало, что наибольший выигрыш для одного цикла публикации карты даёт финальная модификация -- карта проходимости пространства с полным обновлением без блокировок: коэффициент линейной аппроксимации зависимости времени итерации от размера карты уменьшается в $13,5$ раза относительно исходной реализации.
	\item Предложено расширение последней модификации карты проходимости пространства для двумерного пространства.
\end{enumerate}

По результатам работы могут быть даны следующие рекомендации. При построении карты проходимости пространства на БТС, оснащённых камерами и лидарами, проецирование результатов обработки изображения целесообразно выполнять на оценённую по дальномерным данным поверхность движения, а не на фиксированную плоскость: как показано во второй главе, отказ от предположения о плоской поверхности улучшает описание сцены даже по сравнению с алгоритмом, вовсе не использующим визуальные данные. Объединение фрагментов карты, построенных по разным источникам данных, следует выполнять конкурирующим комплексированием, переводя элементы карты, о проходимости которых источники дают противоречивые сведения, в состояние <<неизвестно>>. Наконец, при проектировании программного обеспечения построения карты источники данных целесообразно подключать через единый интерфейс поставщика фрагментов, что позволяет наращивать их число без изменения кода построения карты.

Перспективы дальнейшей разработки темы состоят в следующем. Двумерная реализация карты проходимости пространства с полным обновлением без блокировок требует переноса и валидации на реальной платформе, а также измерения задержки обновления карты при различном числе источников данных в тех же условиях, в которых выполнено профилирование программного комплекса. Это позволит определить условия, при которых она вычислительно эффективнее исходной реализации. Разработанный программный комплекс допускает подключение поставщиков фрагментов карты по данным радаров и сонаров, однако качество получаемой при этом модели в настоящей работе не исследовалось. Отдельным направлением является переход к динамической карте проходимости пространства, хранящей вместе с занятостью элемента оценку скорости наблюдаемого препятствия, что позволило бы отказаться от механизма затухания как основного способа учёта подвижных объектов. Отдельного внимания заслуживает используемая в работе метрика PFC-MSE: стоимости путей в ней вычисляются без учёта кинематической модели движения и габаритов БТС, из-за чего области, недостижимые для БТС, дают вклад в итоговое значение метрики наравне с достижимыми. Учёт кинематической модели движения и габаритов БТС при построении графа сделал бы оценку влияния ошибок карты на алгоритмы планирования пути ближе к тому, как карта используется в действительности. Наконец, представляет интерес оценка предложенных алгоритмов на публичных наборах данных: их применение потребует либо подбора наборов с сопоставимой конфигурацией датчиков, либо построения эталонной разметки проходимости пространства по имеющимся в них аннотациям.
